Positive candidate effects at $D=600$ persist across the three switching
schedules (Figure~\ref{fig:schedule}). Forced hard-minus-raw within-edge
temporal AUROC is $+0.072$ $[0.055,0.089]$ for fixed phase, $+0.072$
$[0.050,0.094]$ for random phase and $+0.070$ $[0.047,0.091]$ for
irregular runs; all three pass secondary Holm6 correction. At $D=300$,
the estimates are $-0.025$, $+0.003$ and $-0.019$, with intervals crossing
zero. The paired duration differences are $+0.097$, $+0.070$ and $+0.089$,
all passing the separate secondary Holm3 family. All six cells have 50
successful, evaluable repeats (300 condition fits in total). The fixed-phase records contain eleven/five
window-centre switches, extending the three/one-switch contrast of Study A.

Gate acceptance changes the effect of executing this candidate. Acceptance
is $2/50$, $0/50$, $3/50$ at $D=300$ and $6/50$, $6/50$, $7/50$
at $D=600$ for fixed, random and irregular schedules. Gated hard aggregation
with static fallback has negative gains throughout, from $-0.214$ to
$-0.198$ (Figure~\ref{fig:gate}). Section~\ref{sec:control-results} examines
how much of that loss comes from the rejection output.

The factor diagnostics in Table~\ref{tab:factors} account for the stronger
forced candidates at the longer duration. Mean per-repeat $JC_{\rm edge}$
is approximately $0.666$, $0.656$ and $0.659$ at $D=600$, compared with
$0.358$, $0.400$ and $0.363$ at $D=300$. State-assignment discrimination
and graph-contrast concordance change together in these comparisons.
